\section{System Overview}
\label{sec:chains}

\subsection{Derivation Pipeline}

\sys takes a program expressed in TVM Tensor Expressions (TE) and the declared domain of its inputs.
As shown in \cref{fig:pipeline}, it first extracts dependent reduction chains over a shared axis (\cref{sec:input-chains}) and then derives a single-pass replacement for each chain.
For each extracted chain, algebraic analysis separates each reduction's context from its input-dependent terms.
The resulting reference states join independently and determine the outputs, establishing a sufficient single-pass construction and a default limit on the maintained state count (\cref{sec:existence}).

\reservewrapspace{26\baselineskip}
\begin{wrapfigure}[26]{r}{0.48\textwidth}
\centering
% Shared vector source for the paper and standalone diagram.
\begin{tikzpicture}[node distance=3mm,
  process/.style={stage,text width=4.5cm,font=\sffamily\small,minimum height=7mm}]
\node[process,draw=origblue] (extract) {Extract dependent reduction\\chains (\S\ref{sec:input-chains})};
\node[process,below=of extract] (judge) {Establish a single-pass form\\State-count upper bound (\S\ref{sec:existence})};
\node[process,below=of judge,draw=auxpurple] (candidate) {Reference tuple from source states\\Synthesise auxiliary states\\(\S\ref{sec:selection}--\ref{sec:synthesis})};
\node[process,below=of candidate] (derive) {Construct join through recovery\\Verify correction formulas\\(\S\ref{sec:recovery}--\ref{sec:rebase})};
\node[process,below=of derive] (check) {Emit and check fused program\\(\S\ref{sec:accuracy}--\ref{sec:impl})};
\node[process,below=of check,draw=fixgreen] (output) {Accept single-pass rewrite};
\foreach \a/\b in {extract/judge,judge/candidate,candidate/derive,derive/check,check/output}
  \draw[flow] (\a) -- (\b);
\draw[flow,draw=auxpurple] (derive.west) -- ++(-0.35,0) |- (candidate.west);
\node[lbl,rotate=90,anchor=south] at ([xshift=-3.5mm]derive.north west) {retry};
\node[lbl,below=2mm of output] {Unaccepted chains remain unchanged};
\end{tikzpicture}

\caption{End-to-end derivation.}
\label{fig:pipeline}
\end{wrapfigure}
Selection starts with the source reduction results as the initial maintained states.
For a source result that is undefined on a legal non-empty part, the procedure seeks a defined representation by extracting the offending context and retaining its coefficients.
Expanding the definitions of these initial states yields their reference tuple and the expressions reconstructing them from it.
This tuple fixes the information that auxiliary-state synthesis must retain.
To discover missing information, the analyser derives how each retained result changes when its local context is replaced by the context of a larger part.
The coefficients of these corrections provide candidate auxiliary states.
The procedure first checks the initial tuple alone, then this tuple augmented with all admissible correction coefficients.
Only if neither tuple succeeds within the state limit does syntax-guided synthesis search combinations of these candidate expressions and, subsequently, reference states.
The search therefore concerns only auxiliary states that, together with the source states, recover the reference tuple (\cref{sec:selection,sec:synthesis}).

For each proposed maintained tuple, verification constructs expressions recovering this entire reference tuple and proves their identities and definedness on legal parts.
This recovery defines a symbolic join: recover the reference states of the two parts, combine them independently, and express the result in maintained states.
Verification establishes that the joined tuple is the state tuple of the union; unsuccessful proofs return to state selection within the search limits.
For execution, context-correction analysis constructs formulas on the maintained states and their local contexts, and verifies them against this symbolic join.
These formulas reuse source centring and scaling where possible to retain its numerical structure; a component without such a formula uses the symbolic join expression (\crefrange{sec:recovery}{sec:rebase}).

The backend emits a TE program that initialises a singleton state tuple for each input element, joins the tuples, and reconstructs the outputs from the maintained states using the source program's scalar dependencies.
Static checks on the reducer and differential numerical measurements in the rebuilt program determine acceptance of this floating-point implementation (\cref{sec:accuracy,sec:impl}).
If no verified construction is found, or emission or numerical checks fail, the source chain is retained.


\subsection{Input Reduction Chains}
\label{sec:input-chains}

The extracted chains contain reductions and intervening computations over a shared axis.
\Cref{fig:grammar} gives the language of reduction bodies; the definitions below specify the dependencies within a chain and its meaning on an input part.

\begin{figure}[htbp]
\small
\begin{align*}
f \;::=\;& x \mid \ctx \mid c \mid N
 \mid f + f \mid f - f \mid f \times f \mid f \div f
 \mid \exp f \mid \log f \mid \sqrt{f} \mid f^{k} \mid |f| \mid \max(f,f) \mid \min(f,f)\\
m \;::=\;& \textstyle\bigoplus_{j} f \mid g(\ctx_1,\dots,\ctx_p)
\end{align*}
\caption{Reduction bodies $f$ and chain members $m$. $x$ ranges over the values a chain reads at position $j$ of the axis, $\ctx$ over earlier members, $c\in\R$ over constants, $k\in\mathbb{Z}$, $N$ is the extent of the axis, and $\oplus\in\{\Sigma,\max,\min\}$.}
\label{fig:grammar}
\end{figure}

\begin{definition}[Reduction chain]
\label{def:chain}
Fix an axis $\Ax=\{0,\dots,N-1\}$ and, for each position $j\in\Ax$, the tuple $x_j$ of values that the chain reads at $j$ from tensors outside it.
A \emph{reduction chain} is a sequence of members $m_1,\dots,m_n$ in producer order.
A member is either a \emph{reduction} $m=\bigoplus_{j\in\Ax} f_m(x_j,\ctx_m)$ with $\oplus\in\{\Sigma,\max,\min\}$, or a \emph{value} $m=g_m(\ctx_m)$.
The \emph{contexts} $\ctx_m$ are the final values of earlier reductions, or scalar computations derived from them, that $m$ reads at an index that does not involve $j$.
They are invariant along the axis, although computing them may require completing earlier reductions.
A chain is \emph{dependent} if some reduction reads a context.
\end{definition}

In the running variance example, the final sum reads the mean as a context.
\sys extracts chains from TE programs by treating reductions over the shared axis as chain members and inlining elementwise tensors read along $j$ into their bodies.
A reduction over another axis, such as $QK^\top$ in attention, stays outside the chain and supplies values $x_j$.

To speak about a part of the axis, we evaluate the chain on that part alone.

\begin{definition}[Semantics on an input part]
\label{def:partial}
For a non-empty $S\subseteq\Ax$, $m(S)$ is the value of member $m$ when every reduction is evaluated over $S$ only and every context is its own value on $S$.
The interpretation of extent reads is fixed for each chain before analysis: reads of $N$ denote $|S|$ when the extent represents the size of the current part, or retain the full-input value $N$ when it is a fixed program parameter.
Both interpretations agree on the full input.
The program computes $m(\Ax)$.
\end{definition}

The declared input domain constrains the element values.
A legal part is any non-empty part of an input satisfying these constraints.
An original member may be undefined on such a part even when the full-input output is defined; a maintained state must be defined on every legal part.
A single pass must produce $m(\Ax)$, wherever the source output is defined, from information gathered part by part.

\subsection{Single-Pass Reducer Contract}

The generated program must initialise a state tuple from each element, join these tuples without revisiting the input, and recover the required outputs.
The following contract states the relation that each join must preserve over the reals.

\begin{definition}[Single-pass reducer]
\label{def:reducer}
A \emph{state map} is a tuple $\st=(\st_1,\dots,\st_k)$ of functions of non-empty parts.
For fixed non-reduction indices, each function is scalar-valued, so $\st(S)$ is a state tuple with components $\st_i(S)$; its size $k$ is fixed for a given reducer, independently of the reduction length $N$.
A \emph{single-pass reducer} for a chain is a tuple $(\leaf,\Mrg,\epi)$ of an initialisation function $\leaf(x_j)=\st(\{j\})$, a binary join $\Mrg$ and an output function $\epi$ such that, for all disjoint non-empty parts $\pa$ and $\pb$,
\begin{equation}
\label{eq:invariant}
\Mrg(\st(\pa),\st(\pb))=\st(\pa\dunion\pb),\qquad
\epi(\st(\Ax))=\bigl(m(\Ax)\bigr)_{m\text{ required}}.
\end{equation}
The \emph{required} members are those the rest of the program reads.
\end{definition}

The initialisation function produces a singleton state tuple for each input element; the generated reduction starts from the first tuple and folds the remaining tuples.
No numerical state tuple represents an empty part: execution skips empty parts, and the input axis must be non-empty.
The join invariant of \cref{eq:invariant} connects the maintained state tuple to the original program.

\begin{lemma}[Laws from the invariant]
\label{lem:laws}
If the invariant holds, $\Mrg$ is commutative and associative when applied to state tuples of pairwise disjoint non-empty parts.
For every partition of $\Ax$ into non-empty parts and every binary join tree over them, joining the leaves yields $\st(\Ax)$.
The serial fold initialised from its first singleton state tuple has the same property.
\end{lemma}

\begin{proof}
$\Mrg(\st(\pa),\st(\pb))=\st(\pa\dunion\pb)=\st(\pb\dunion\pa)=\Mrg(\st(\pb),\st(\pa))$, and associativity follows in the same way from that of $\dunion$.
The rest is induction on the join tree and on the length of the fold.
\end{proof}

Each non-empty part can therefore compute its state tuple independently.
A tree of joins combines these tuples without another traversal of the parts.
This parallel structure follows from the invariant, independently of a particular schedule.

In the vocabulary of program calculation, $S\mapsto\st(S)$ is a homomorphism from finite non-empty bags under $\dunion$ to the state tuples under $\Mrg$: equivalently, a list homomorphism~\citep{bird87,gibbons96} invariant under permutation of its input.
In the vocabulary of databases, a chain with a context-independent reference construction is an \emph{algebraic} aggregate, computable from a fixed number of \emph{distributive} ones~\citep{datacube}.
Empty parts are handled by control flow, outside the algebraic join contract.

\paragraph{Problem.}
Given a dependent chain, establish a sufficient condition for a single-pass reducer, construct one, prove \cref{eq:invariant}, and accept it only if in floating point it introduces no additional non-finite result on the checked inputs and its measured error stays within the stated acceptance threshold relative to the source program.
